% ----------------------------------------------------------------
\subsection{La griglia temporale}\label{sec:griglia}

Trattiamo due proprietà: il passo con cui i grani si succedono (density) e la regolarità di
quel passo (distribution).

% ----------------------------------------------------------------
\subsubsection{La density}

Il passo della griglia è l'intervallo fra due onset successivi,
l'\emph{inter-onset time}, IOT\cite{Truax1988}. 
Come anticipato precedentemente \S\ref{sec:c-e}, lo stato di default
usa il $\texttt{fill\_factor} = 2$\footnote{Il \emph{fill factor} di Roads è
il prodotto della densità per la durata del grano, $\mathrm{FF} =
\mathit{density}\cdot\texttt{grain.duration}$: il numero medio di grani che
coprono un istante --- a $\mathrm{FF}=1$ adiacenti (\emph{covered}), sotto 1
\emph{sparse}, sopra 1 \emph{packed}~\cite[pp.~105--106]{Roads2001}. Coincide
col fattore di overlap: il default 2 copre ogni istante con due grani.}, da cui
la densità discende come funzione della durata del grano,
$\texttt{density} = \texttt{fill\_factor}\,/\,\texttt{grain.duration}$.
Il riempimento resta così invariante rispetto alla durata: grani più
brevi si infittiscono, grani più lunghi si diradano, e il numero medio
di grani sovrapposti non cambia.

\begin{figure}[h]
    \includegraphics[width=\linewidth]{examples/distribution/distribution_map}\\
    \captionof{figure}{La \textsc{map} del rendering: ascisse = tempo del
    brano, ordinate = posizione di lettura nel buffer (costante, per il
    congelamento); ogni segmento un grano.}
    \label{fig:griglia-map}
\end{figure}


Questa invarianza è comoda finché densità e durata del grano si muovono
insieme, ma diventa un vincolo quando devono variare in modo
indipendente --- come nei frammenti più estesi che seguono. In quel
caso si dichiara \texttt{density} in forma esplicita, in grani al
secondo, e il \texttt{fill\_factor} cede il posto. È la forma che adotto
da qui in avanti: \texttt{density} esplicita fissa il passo medio della
griglia $\overline{\mathrm{IOT}} = 1/\mathit{density}$, ed è su questo
passo che interviene la deviazione successiva.

\subsubsection{La distribution}

Fissato il passo medio, \texttt{distribution} ne governa la regolarità.
Il parametro interpola fra due regimi:

\begin{equation}\label{eq:iot}
  \mathrm{IOT}(d) = (1-d)\,\overline{\mathrm{IOT}} + d\,\mathcal{U}\!\left(0, 2\overline{\mathrm{IOT}}\right), \quad d \in [0,1].
\end{equation}

A \texttt{distribution} $=0$ ogni intervallo vale
$\overline{\mathrm{IOT}}$: l'emissione è \emph{sincrona}, la griglia
metronomica. A \texttt{distribution} $=1$ l'intervallo è campionato da
una distribuzione uniforme su $[0,\,2\overline{\mathrm{IOT}}]$:
l'emissione è \emph{asincrona}.\footnote{\emph{Sincrono} e
\emph{asincrono} sono qui nel senso di Truax~\cite{Truax1988} --- il
grado di regolarità degli inter-onset --- non in quello, pure corrente
nel lessico granulare, di sintesi \emph{pitch-synchronous} o
\emph{period-synchronous}~\cite{DePoliPiccialli1988}, dove la sincronia
è con il periodo della forma d'onda e serve a fissare l'altezza.} I
valori intermedi pesano linearmente i due regimi. Poiché il valore
atteso della distribuzione uniforme è $\overline{\mathrm{IOT}}$, la
densità media si conserva per ogni $d$: ciò che cambia lungo la
transizione è la regolarità degli istanti, non il loro numero.

    \lstinputlisting[
      linerange={5,10,11,16,21,22,27},
      xleftmargin=-1em,
      language=yaml,
      basicstyle=\ttfamily\fontsize{7.5}{9}\selectfont,
      label=lst:distribution,
      caption={\texttt{distribution} come envelope da 0 a 1
      (\texttt{ex1\_distribution.yml}).}
    ]{examples/distribution/distribution.yml}

Su una densità abbastanza rada da risolvere i singoli onset, la
\textsc{map} mostra la transizione direttamente
(Fig.~\ref{fig:griglia-map}): in ascissa il tempo del brano, in
ordinata la posizione di lettura nel buffer --- qui costante, per il
congelamento --- e ogni segmento un grano; da sinistra a destra la
spaziatura passa da uniforme a dispersa, a parità di numero medio di
grani.

L'effetto non è solo ritmico. Quando la densità è alta e il passo
regolare, la griglia è essa stessa periodica, e un treno periodico di
grani ha uno spettro a righe, spaziate a multipli della densità ---
le bande laterali che a \S\ref{sec:c-e} si elidevano nella
sovrapposizione, e che qui non si elidono, perché la lettura è ferma e
i grani sono identici~\cite{Roads2001}. Irregolarizzare gli intervalli
disperde quella periodicità: lo spettro a righe si spalma.

Rompere questa periodicità è possibile per più vie, distinte dall'asse
su cui agiscono. \texttt{distribution} agisce sul \emph{tempo} della
griglia: irregolarizza gli istanti di emissione. Le stesse righe si
disperdono agendo invece sul \emph{contenuto} dei grani --- il punto del
buffer che ciascuno legge, la sua ampiezza --- ed è la deviazione per
grano di \S\ref{sec:deviazione}. Nel congelamento le due vie sono
complementari: \texttt{distribution} lascia i grani identici nel
contenuto ma irregolari nel tempo; resta da sciogliere l'identità di
contenuto.